\chapter{Data model}
\label{ch:data-model}

The database holds eight application tables in the \code{public} schema and one log table in the
\code{private} schema. Every table references \code{profiles}, which has one row per account and itself
references the Supabase Auth user. \Cref{fig:er} shows the tables, their columns and the foreign keys
between them. Every foreign key to a user cascades on delete, so removing an Auth user removes everything
that belongs to that user (\cref{sec:flow-delete}).

\begin{figure}[p]
  \centering
  % Entity-relationship diagram of the database (Figure: entity-relationship).
% Each box lists a table's columns; PK marks the primary key and FK a foreign key. Lines are foreign
% keys, labelled with the referencing columns; 1, n and 0..1 give the cardinality at each end.
\begin{tikzpicture}[
    k/.style={font=\sffamily\tiny, text=muted, inner sep=1pt},
    kl/.style={font=\sffamily\tiny, text=ink, inner sep=1pt},
  ]
  \newcommand{\pk}{\hfill{\sffamily\tiny PK}}
  \newcommand{\fk}{\hfill{\sffamily\tiny FK}}
  \newcommand{\pkfk}{\hfill{\sffamily\tiny PK, FK}}

  % ---------------------------------------------------------------- column B
  \node[entity=30mm, anchor=north] (users) at (7.95,0) {\textbf{auth.users}\nodepart{two}
    id\pk\\ email\\ {\sffamily\tiny managed by Supabase Auth}};
  \node[entity=30mm, anchor=north, draw=accent, line width=0.8pt] (profiles) at (7.95,-2.5) {\textbf{profiles}\nodepart{two}
    id\pkfk\\ first\_name\\ last\_name\\ gender\\ show\_gender\\ intro\\ interests\\ photo\_path\\
    photo\_paths\\ is\_adult\\ verified\\ status\\ created\_at};
  \node[entity=30mm, anchor=north] (members) at (7.95,-8.6) {\textbf{event\_members}\nodepart{two}
    event\_id\pkfk\\ user\_id\pkfk\\ joined\_at};
  \node[privateentity=30mm, anchor=north] (posts) at (7.95,-10.9) {\textbf{event\_posts}\nodepart{two}
    user\_id\fk\\ posted\_at};

  % ---------------------------------------------------------------- column A
  \node[entity=30mm, anchor=north, dashed] (objects) at (1.75,0) {\textbf{storage.objects}\nodepart{two}
    bucket\_id = 'photos'\\ name = <owner>/<file>};
  \node[entity=30mm, anchor=north] (availability) at (1.75,-2.5) {\textbf{availability}\nodepart{two}
    user\_id\pkfk\\ location\\ accuracy\_m\\ note\\ state\\ started\_at\\ expires\_at};
  \node[privateentity=30mm, anchor=north] (changes) at (1.75,-6.2) {\textbf{location\_changes}\nodepart{two}
    user\_id\fk\\ changed\_at};
  \node[entity=30mm, anchor=north] (blocks) at (1.75,-8.4) {\textbf{blocks}\nodepart{two}
    blocker\pkfk\\ blocked\pkfk\\ created\_at};
  \node[entity=30mm, anchor=north] (reports) at (1.75,-10.8) {\textbf{reports}\nodepart{two}
    id\pk\\ reporter\fk\\ reported\fk\\ request\_id\fk\\ reason\\ details\\ created\_at\\ reviewed\_at\\ action};

  % ---------------------------------------------------------------- column C
  \node[entity=30mm, anchor=north] (requests) at (14.1,-0.6) {\textbf{chat\_requests}\nodepart{two}
    id\pk\\ from\_user\fk\\ to\_user\fk\\ status\\ created\_at\\ expires\_at\\ responded\_at\\ ended\_at};
  \node[entity=30mm, anchor=north] (messages) at (14.1,-4.9) {\textbf{messages}\nodepart{two}
    id\pk\\ request\_id\fk\\ event\_id\fk\\ sender\fk\\ sender\_name\\ body\\ created\_at};
  \node[entity=30mm, anchor=north] (events) at (14.1,-9.0) {\textbf{events}\nodepart{two}
    id\pk\\ host\fk\\ emoji\\ title\\ description\\ location\\ starts\_at\\ created\_at};

  % ---------------------------------------------------------------- relationships
  % profiles - auth.users
  \draw[relation] (users.south) -- node[kl, right=2pt] {id} (profiles.north);
  \node[k, left=2pt] at ($(users.south)+(0,-0.15)$) {1};
  \node[k, left=2pt] at ($(profiles.north)+(0,0.15)$) {1};
  % profiles - storage.objects: by folder name, not a foreign key
  \draw[relation, dashed] (objects.east |- 0,-0.7) -- (5.2,-0.7) |- (profiles.west |- 0,-2.9);
  \node[kl, anchor=west] at (5.25,-1.75) {folder = id};
  % profiles - availability
  \draw[relation] (profiles.west |- 0,-3.6) -- node[kl, above=1pt] {user\_id} (availability.east |- 0,-3.6);
  \node[k, below=1pt] at ($(profiles.west |- 0,-3.6)+(-0.18,0)$) {1};
  \node[k, below=1pt] at ($(availability.east |- 0,-3.6)+(0.3,0)$) {0..1};
  % profiles - location_changes
  \draw[relation] (profiles.west |- 0,-6.85) -- node[kl, above=1pt] {user\_id} (changes.east |- 0,-6.85);
  \node[k, below=1pt] at ($(profiles.west |- 0,-6.85)+(-0.18,0)$) {1};
  \node[k, below=1pt] at ($(changes.east |- 0,-6.85)+(0.18,0)$) {n};
  % profiles - blocks and reports, through one trunk
  \draw[relation] (profiles.west |- 0,-7.3) -- (5.1,-7.3) |- (blocks.east |- 0,-9.3);
  \draw[relation] (5.1,-9.3) |- (reports.east |- 0,-11.6);
  \node[kl, above=1pt] at ($(blocks.east |- 0,-9.3)!0.5!(5.1,-9.3)$) {blocker, blocked};
  \node[kl, above=1pt] at ($(reports.east |- 0,-11.6)!0.5!(5.1,-11.6)$) {reporter, reported};
  \node[k, below=1pt] at ($(profiles.west |- 0,-7.3)+(-0.18,0)$) {1};
  \node[k, below=1pt] at ($(blocks.east |- 0,-9.3)+(0.18,0)$) {n};
  \node[k, below=1pt] at ($(reports.east |- 0,-11.6)+(0.18,0)$) {n};
  % profiles - event_posts, branching from the same trunk
  \draw[relation] (5.1,-11.35) -- node[kl, above=1pt] {user\_id} (posts.west |- 0,-11.35);
  \node[k, below=1pt] at ($(posts.west |- 0,-11.35)+(-0.18,0)$) {n};
  % profiles - chat_requests
  \draw[relation] (profiles.east |- 0,-3.1) -- node[kl, above=1pt] {from\_user, to\_user} (requests.west |- 0,-3.1);
  \node[k, below=1pt] at ($(profiles.east |- 0,-3.1)+(0.18,0)$) {1};
  \node[k, below=1pt] at ($(requests.west |- 0,-3.1)+(-0.18,0)$) {n};
  % profiles - messages
  \draw[relation] (profiles.east |- 0,-6.3) -- node[kl, above=1pt] {sender} (messages.west |- 0,-6.3);
  \node[k, below=1pt] at ($(profiles.east |- 0,-6.3)+(0.18,0)$) {1};
  \node[k, below=1pt] at ($(messages.west |- 0,-6.3)+(-0.18,0)$) {n};
  % profiles - events
  \draw[relation] (profiles.east |- 0,-7.3) -- (11.0,-7.3) |- (events.west |- 0,-9.4);
  \node[kl, above=1pt] at ($(11.0,-9.4)!0.5!(events.west |- 0,-9.4)$) {host};
  \node[k, below=1pt] at ($(profiles.east |- 0,-7.3)+(0.18,0)$) {1};
  \node[k, below=1pt] at ($(events.west |- 0,-9.4)+(-0.18,0)$) {n};
  % profiles - event_members
  \draw[relation] (members.north) -- node[kl, right=2pt] {user\_id} (profiles.south);
  \node[k, left=2pt] at ($(profiles.south)+(0,-0.15)$) {1};
  \node[k, left=2pt] at ($(members.north)+(0,0.15)$) {n};
  % events - event_members
  \draw[relation] (members.east |- 0,-9.95) -- node[kl, below=1pt] {event\_id} (events.west |- 0,-9.95);
  \node[k, above=1pt] at ($(members.east |- 0,-9.95)+(0.18,0)$) {n};
  \node[k, above=1pt] at ($(events.west |- 0,-9.95)+(-0.18,0)$) {1};
  % chat_requests - messages, events - messages
  \draw[relation] (requests.south) -- node[kl, right=2pt] {request\_id} (messages.north);
  \node[k, left=2pt] at ($(requests.south)+(0,-0.15)$) {0..1};
  \node[k, left=2pt] at ($(messages.north)+(0,0.15)$) {n};
  \draw[relation] (messages.south) -- node[kl, right=2pt] {event\_id} (events.north);
  \node[k, left=2pt] at ($(messages.south)+(0,-0.15)$) {n};
  \node[k, left=2pt] at ($(events.north)+(0,0.15)$) {0..1};
  % chat_requests - reports
  \draw[relation] (reports.east |- 0,-14.1) -- (15.95,-14.1) |- (requests.east |- 0,-1.15);
  \node[kl, above=1pt] at (9.6,-14.1) {request\_id (set to null when the request is deleted)};
  \node[k, above=1pt] at ($(reports.east |- 0,-14.1)+(0.18,0)$) {n};
  \node[k, above=1pt] at ($(requests.east |- 0,-1.15)+(0.12,0.02)$) {0..1};
\end{tikzpicture}

  \caption{Entity-relationship diagram. Each box lists a table's columns. Lines are foreign keys,
    labelled with the referencing columns, with the cardinality at each end. The dashed box is the
    Storage table that holds photos; it is linked to \code{profiles} by folder name, not by a foreign
    key. The dashed table is in the \code{private} schema.}
  \label{fig:er}
\end{figure}

\section{Tables}
\label{sec:tables}

\subsection{Profiles}

\code{profiles} has one row per account, keyed by the Auth user's ID. It holds what the user entered
during profile setup: first and last name (1 to 30 characters each), gender (one of \code{woman},
\code{man}, \code{non-binary} and \code{prefer-not}), whether to show the gender, an introduction (1 to
300 characters) and 3 to 5 interests. \code{photo_path} is the storage path of the profile picture and
\code{photo_paths} the paths of up to six further photos. Check constraints require every path to name a
file directly inside the owner's own folder (\cref{sec:photo-privacy}).

Three columns are controlled by the server. \code{is_adult} records the 18+ confirmation and must be
true when the row is inserted. \code{verified} is reserved for the selfie check; clients cannot write it.
\code{status} is \code{active}, \code{hidden} or \code{suspended}; reports set it to \code{hidden}
automatically (\cref{sec:blocking-reporting}), and only an \code{active} profile can go live, appear in
lists, create or join events, or receive a request.

Other users never read this table. What they may see of a profile is built by the helper
\code{private.person_card} (\cref{sec:person-card}).

\subsection{Availability}

\code{availability} holds one row for each user who is live right now, and nothing else. The row stores
the position as a PostGIS \code{geography(point, 4326)}, the accuracy the phone reported, the live note
(up to 60 characters), a state of \code{open} or \code{chatting}, and the start and end of the session.
Rows are deleted when the user goes offline and by the minute job when the session ends, so the table
stays small however many accounts exist, and its spatial index stays fast. No client can read the table;
the position is used only inside database functions and never returned.

The state column distinguishes a user who can be found (\code{open}) from one who is in a chat
(\code{chatting}) and is therefore left out of nearby lists. \Cref{fig:live-states} shows how the row
moves between these states. \code{go_green} inserts the row as \code{open}, or as \code{chatting} if the
user still has an accepted chat, which happens when a user goes offline during a chat and then goes live
again. \code{respond} sets both users to \code{chatting} when a request is accepted. Every way a chat can
end sets both users back to \code{open} if they are still live: \code{end_chat}, \code{cancel_request}
after an accept, \code{block_user} when it ends a chat between the two, and \code{delete_account} for the
partner of the deleted user. \code{go_offline} deletes the row, and \code{expire_stale} deletes it at the
session's end time.

\begin{figure}[ht]
  \centering
  % Live state of a user: the availability row and its state column (Figure: live states).
\begin{tikzpicture}
  \node[endstate, minimum width=24mm] (off) at (0,0) {not live\\[-1pt]{\scriptsize\mdseries no row}};
  \node[state, minimum width=24mm] (open) at (6.6,0) {open\\[-1pt]{\scriptsize\mdseries on nearby lists}};
  \node[state, minimum width=24mm] (chatting) at (13.2,0) {chatting\\[-1pt]{\scriptsize\mdseries hidden from lists}};

  \draw[transition] ($(off.east)+(0,0.2)$) -- node[translabel, above=2pt, fill=none] {goes live} ($(open.west)+(0,0.2)$);
  \draw[transition] ($(open.west)+(0,-0.2)$) -- node[translabel, below=2pt, fill=none] {goes offline, or time is up} ($(off.east)+(0,-0.2)$);
  \draw[transition] ($(open.east)+(0,0.2)$) -- node[translabel, above=2pt, fill=none] {a request is accepted} ($(chatting.west)+(0,0.2)$);
  \draw[transition] ($(chatting.west)+(0,-0.2)$) -- node[translabel, below=2pt, fill=none] {the chat ends} ($(open.east)+(0,-0.2)$);
  \draw[transition] (chatting.north) -- ++(0,0.7) -| node[translabel, pos=0.25, above=1pt, fill=none] {goes offline, or time is up} (off.north);
  \draw[transition] (off.south) -- ++(0,-1.0) -| node[translabel, pos=0.25, below=1pt, fill=none] {goes live while still in an accepted chat} (chatting.south);
\end{tikzpicture}

  \caption{Live state of one user. ``Not live'' means the user has no \code{availability} row. Any chat
    ending returns both users to \code{open} if they are still live.}
  \label{fig:live-states}
\end{figure}

\subsection{Requests and chats}

\code{chat_requests} records every say-hi request. A request names its sender (\code{from_user}) and
recipient (\code{to_user}), which must differ, and has a status. It expires five minutes after it is
created unless it is answered. When the recipient accepts, the same row becomes the chat: its ID is the
chat's ID, its status stays \code{accepted} while the chat lasts, and \code{ended_at} records when it
ended. \code{responded_at} records when the recipient answered. A partial unique index allows at most one
\code{pending} request per sender.

\Cref{fig:request-states} shows the lifecycle of the status column and \cref{tab:request-transitions}
lists every transition with the code that performs it.

\begin{figure}[ht]
  \centering
  % Lifecycle of a say-hi request: the status column of chat_requests (Figure: request states).
% Labels are short; the table of request transitions gives every trigger.
\begin{tikzpicture}
  \node[startdot] (start) at (0,0) {};
  \node[state] (pending) at (2.3,0) {pending};
  \node[state] (accepted) at (8.4,0) {accepted};
  \node[endstate] (met) at (13.2,0) {met};
  \node[endstate] (declined) at (0.6,-3.4) {declined};
  \node[endstate] (expired) at (3.9,-3.4) {expired};
  \node[endstate] (cancelled) at (8.4,-3.4) {cancelled};

  \draw[transition] (start) -- (pending);
  \node[translabel, fill=none, above=3pt] at ($(start)!0.5!(pending.west)$) {T1};
  \draw[transition] (pending) -- node[translabel, above=2pt, fill=none] {T2\enskip accept} (accepted);
  \draw[transition] (accepted) -- node[translabel, above=2pt, fill=none] {T7\enskip ended, met} (met);
  \draw[transition] ($(pending.south)+(-0.5,0)$) -- node[translabel, left=3pt, fill=none] {T3\enskip decline} (declined.north);
  \draw[transition] ($(pending.south)+(0.4,0)$) -- node[translabel, right=3pt, fill=none] {T4\enskip expired} (expired.north);
  \draw[transition] ($(pending.east)+(0,-0.22)$) -- (6.0,-0.22) |- (cancelled.west);
  \node[translabel, fill=none, align=left, anchor=west] at (6.1,-1.75) {T5\enskip withdrawn\\\hphantom{T5\enskip}or blocked};
  \draw[transition] (accepted.south) -- node[translabel, right=3pt, fill=none, align=left] {T6\enskip ended, not met;\\\hphantom{T6\enskip}withdrawn; blocked;\\\hphantom{T6\enskip}or stale} (cancelled.north);
\end{tikzpicture}

  \caption{Lifecycle of a request. Grey states are final. The labels refer to
    \cref{tab:request-transitions}.}
  \label{fig:request-states}
\end{figure}

\begin{table}[ht]
  \centering
  \caption{Request transitions.}
  \label{tab:request-transitions}
  \tablefont
  \begin{tabularx}{\textwidth}{@{}L{8mm}L{44mm}Y@{}}
    \toprule
    & \textsf{\textbf{Transition}} & \textsf{\textbf{Performed by}} \\
    \midrule
    T1 & created as \code{pending} & \code{send_request}, called by the sender \\
    T2 & \code{pending} to \code{accepted} & \code{respond} with \code{accept = true}, called by the recipient \\
    T3 & \code{pending} to \code{declined} & \code{respond} with \code{accept = false}, called by the recipient \\
    T4 & \code{pending} to \code{expired} & \code{expire_stale}, once five minutes have passed or the sender is no longer live; \code{send_request} also expires the sender's own request whose time is up before checking for one already waiting \\
    T5 & \code{pending} to \code{cancelled} & \code{cancel_request} (the sender withdraws it); \code{go_offline} (the sender goes offline); \code{respond} when the sender accepts someone else's request; \code{block_user} by either person \\
    T6 & \code{accepted} to \code{cancelled} & \code{end_chat} with outcome \code{cancelled}; \code{cancel_request} when the cancel arrives just after an accept; \code{block_user} by either person; \code{expire_stale} once neither person is live \\
    T7 & \code{accepted} to \code{met} & \code{end_chat} with outcome \code{met} \\
    \bottomrule
  \end{tabularx}
\end{table}

The sender never learns why a request did not lead to a chat. The client treats every final status other
than \code{accepted} as \ui{Not this time} (\cref{sec:flow-sayhi}).

\subsection{Messages}

\code{messages} holds the messages of both kinds of chat. A check constraint requires exactly one parent:
\code{request_id} for a one-to-one chat or \code{event_id} for an event's group chat. The body is 1 to
1000 characters. \code{sender_name} stores the sender's display name (first name and last initial) at the
time of sending, because a message delivered over Realtime must be renderable on its own: the recipient
cannot read the sender's profile row to look the name up.

\subsection{Events and members}

\code{events} holds the public events on the map. Each has a host, a category, a title (1 to 60
characters), a description (up to 200 characters), a position and a start time. The position is a place
the host chose on the map and is, by design, the one location in the system that other users are meant
to see. The \code{emoji} column stores the category as one of twelve emoji values; the client maps each
value to an icon and a name (\cref{tab:event-categories}). \code{event_members} records who has joined
which event. The host is added as a member when the event is created, and only the host can delete the
event (\code{delete_event}); its members and messages are deleted with it.

\subsection{Blocks and reports}

\code{blocks} records that one user has blocked another; the pair is the primary key and a user cannot
block themselves. Although each row has a direction, the server applies a block in both directions
(\cref{sec:blocking-reporting}). \code{reports} records a report: who reported whom, an optional request
the report relates to, one of six reasons (\code{inappropriate-photo}, \code{harassment}, \code{unsafe},
\code{spam-or-fake}, \code{under-18}, \code{other}) and up to 500 characters of details. The columns
\code{reviewed_at} and \code{action} exist for a later human review; nothing in the repository writes
them.

\subsection{Location changes}

\code{private.location_changes} records the time at which each user started a live session or moved more
than 150 metres while live. It stores no position. \code{go_green} counts the rows from the last ten
minutes to limit how often a location may change, and the minute job deletes rows older than that
(\cref{sec:relocation-limits}).

\subsection{Event posts}

\code{private.event_posts} records the time each event was created, and nothing else. \code{create_event}
counts the rows from the last 24 hours for its five-a-day limit, so deleting an event does not give its
place back; the minute job deletes rows older than a day.

\subsection{Photos in Storage}

Photo files are objects in the Storage bucket \code{photos}, which is private, accepts JPEG, PNG and
WebP, and limits each file to 5\,MiB. The app names every object
\code{<user-id>/<slot>-<milliseconds>.jpg}, where the slot is \code{profile} for the profile picture and
\code{photo0} to \code{photo5} for the others. The profile picture also has a 320-pixel copy at the same
path with the suffix \code{-small.jpg}, used for small avatars.

\section{Indexes}

\Cref{tab:indexes} lists the indexes beyond the primary keys and what each one serves.

\begin{table}[ht]
  \centering
  \caption{Secondary indexes.}
  \label{tab:indexes}
  \tablefont
  \begin{tabularx}{\textwidth}{@{}L{78mm}Y@{}}
    \toprule
    \textsf{\textbf{Index}} & \textsf{\textbf{Serves}} \\
    \midrule
    \code{availability (location)}, GiST & Distance queries in \code{nearby} and \code{near_each_other} \\
    \code{availability (expires_at)} & The minute job's deletion of ended sessions \\
    \code{chat_requests (from_user)} where \code{pending}, unique & At most one waiting request per sender \\
    \code{chat_requests (to_user)} & Looking up a recipient's waiting requests \\
    \code{chat_requests (from_user, created_at)} & The rate limit on sending requests \\
    \code{chat_requests (expires_at)} where \code{pending} & Expiring unanswered requests \\
    \code{chat_requests (from_user, to_user)} where \code{accepted} & One-chat checks, \code{my_status} and stale-chat cleanup \\
    \code{events (location)}, GiST & Radius queries in \code{events_near} \\
    \code{events (host)}, \code{events (starts_at)} & Rate limit on creating events; time filter \\
    \code{event_members (user_id)} & Membership lookups by user \\
    \code{messages (request_id, created_at)} & One-to-one chat history \\
    \code{messages (event_id, created_at)} & Group chat history \\
    \code{messages (sender, created_at)} & The rate limit on sending messages \\
    \code{blocks (blocked)} & Block checks in the reverse direction \\
    \code{reports (reported)}, \code{(reporter)}, \code{(request_id)} & Counting reporters; cascades \\
    \code{private.location_changes (user_id, changed_at)} & The relocation limit in \code{go_green} \\
    \code{private.event_posts (user_id, posted_at)} & The event limit in \code{create_event} \\
    \bottomrule
  \end{tabularx}
\end{table}

\section{Data the client works with}
\label{sec:client-types}

The client does not mirror the tables. It works with a small set of shapes defined in
\code{src/data/types.ts}, built by the API layer from what the server returns (\cref{tab:client-types}).

\begin{table}[ht]
  \centering
  \caption{Client data shapes and their sources.}
  \label{tab:client-types}
  \tablefont
  \begin{tabularx}{\textwidth}{@{}L{32mm}L{48mm}Y@{}}
    \toprule
    \textsf{\textbf{Type}} & \textsf{\textbf{Built from}} & \textsf{\textbf{Holds}} \\
    \midrule
    \code{Profile} & the user's own \code{profiles} row & All profile fields, with photos as signed URLs \\
    \code{NearbyPerson} & a person card & What another user may see; photo URLs signed on demand \\
    \code{Availability} & \code{go_green} or \code{my_status} & Duration, note, server start and end times \\
    \code{OutgoingRequest} & \code{send_request} or \code{my_status} & Request ID, recipient, end time, status \code{pending}, \code{accepted} or \code{not-this-time} \\
    \code{IncomingRequest} & a Realtime insert or a query, then \code{request_person} & Request ID, sender's card, end time \\
    \code{ChatSession} & an accepted request & Chat ID, partner, status \code{active} or \code{other-cancelled}, messages \\
    \code{Message} & a \code{messages} row & ID, sender (\code{me} for one's own), sender name, text, time \\
    \code{MapEvent} & an event card & Event fields, host name and photo, member count, joined flag, loaded messages \\
    \bottomrule
  \end{tabularx}
\end{table}

\subsection{The person card}
\label{sec:person-card}

A \term{person card} is the only form in which one user's profile reaches another user. The helper
\code{private.person_card(target)} builds it as JSON with these fields: ID, first name, last initial,
gender (only if the owner chose to show it and did not choose \code{prefer-not}), introduction,
interests, the storage paths of the profile picture and other photos, the live note, and the
\code{verified} flag. It contains no last name, no location and no distance. Every function that returns
a full profile (\code{nearby}, \code{request_person}, \code{my_status}) returns person cards, so this helper
is the single definition of what a profile reveals. Two places return less: \code{blocked_list} returns
only ID, first name and last initial, and an event card's host adds only the profile picture's path to
those.

\subsection{Event categories}

The category of an event is stored as an emoji value in \code{events.emoji}. The app shows an icon and a
name instead; the table \code{EVENT_CATEGORIES} in \code{src/data/constants.ts} maps one to the other
(\cref{tab:event-categories}). A value missing from the table, for example one written by a newer app,
falls back to a calendar icon named \ui{Event}. The database default is the first value, Coffee.

\begin{table}[ht]
  \centering
  \caption{Event categories. Stored values are given as Unicode code points.}
  \label{tab:event-categories}
  \tablefont
  \begin{tabular}{@{}lll@{\hspace{12mm}}lll@{}}
    \toprule
    \textsf{\textbf{Name}} & \textsf{\textbf{Icon}} & \textsf{\textbf{Value}} & \textsf{\textbf{Name}} & \textsf{\textbf{Icon}} & \textsf{\textbf{Value}} \\
    \midrule
    Coffee & \code{cafe} & U+2615 & Sports & \code{basketball} & U+1F3C0 \\
    Food & \code{pizza} & U+1F355 & Music & \code{musical-notes} & U+1F3B8 \\
    Games & \code{dice} & U+1F3B2 & Walk & \code{walk} & U+1F6B6 \\
    Study & \code{library} & U+1F4DA & Yoga & \code{flower} & U+1F9D8 \\
    Coding & \code{laptop} & U+1F4BB & Movies & \code{film} & U+1F3AC \\
    Art & \code{color-palette} & U+1F3A8 & Drinks & \code{beer} & U+1F37B \\
    \bottomrule
  \end{tabular}
\end{table}
